% Glossary terms from ch16: lines of the form \item[Term] Definition.
\item[Sampling-based planning] Planning in a continuous space by drawing random configurations, keeping the collision-free ones and connecting them by short collision-free segments, so that memory grows with the number of samples rather than with the size or dimension of the space.
\item[Curse of dimensionality] The growth of a uniform grid as $N^{d}$ with the resolution $N$ per axis and the dimension $d$, which makes grid search impossible in large 3D volumes and in state spaces that include velocities.
\item[Rapidly-exploring random tree (RRT)] A tree rooted at the start configuration that grows one vertex per random sample by moving at most $\eta$ from the nearest vertex towards the sample when that segment is collision-free; the first vertex in the goal region yields a path.
\item[Voronoi bias] The property that a vertex of the tree is selected for extension with probability equal to the relative volume of its Voronoi region, and extended when the step towards the sample is free, so that vertices facing large unexplored regions grow the tree most often.
\item[Step size] The largest length $\eta$ of a new edge: \textsc{Steer} moves from the nearest vertex towards the sample by at most $\eta$.
\item[Goal bias] The probability $p_{\mathrm{goal}}$ with which \textsc{Sample} returns the goal instead of a uniform configuration; small values ($0.05$--$0.1$) speed the search up, large values make the tree greedy and trapped.
\item[Goal region] The ball of radius $r_{\mathrm{goal}}$ around the goal point; the search ends when a vertex enters it, since continuous samples never hit a point exactly.
\item[Collision checking by subdivision] Testing points spaced at most $\delta$ apart along a segment against the obstacles inflated by $r+\delta/2$; every accepted segment then has clearance at least the robot radius $r$.
\item[Probabilistic completeness] The guarantee that, when a solution with positive clearance exists, the probability of having found one after $n$ iterations tends to one as $n\to\infty$; it says nothing when no solution exists.
\item[Narrow-passage problem] The growth of the expected number of iterations like $(1/w)^{d}$ for a passage of width $w$, because uniform samples rarely fall into it.
\item[RRT-Connect] A bidirectional RRT with one tree from the start and one from the goal, in which the other tree greedily repeats \textsc{Extend} towards every new vertex (\textsc{Connect}) and the trees swap roles every iteration.
\item[Kinodynamic RRT] RRT over states $(\pos,\vel)$ in which \textsc{Steer} is replaced by forward simulation of sampled controls and \textsc{Nearest} uses a weighted state-space metric; the result is a trajectory that obeys the dynamics.
\item[kd-tree] A binary tree that splits points along coordinate axes and answers nearest-neighbour queries in $\bigO{\log n}$ expected time in low dimension, against $\bigO{n}$ for a linear scan.
\item[Shortcutting] Path post-processing that replaces the part of a path between two points by the straight segment when that segment is collision-free; it never increases the length but cannot carry the path across an obstacle.
\item[Probabilistic roadmap (PRM)] The multi-query counterpart of RRT: sample free configurations, connect each to its nearest neighbours by checked segments, and answer queries by graph search on the roadmap.
